%% kapa: exam=1 round=22 year=2011 period=1 form=A subject=civil_law subject_ko=민법 seq=1 number=1
%% source: Q-Net 감정평가사 기출문제 (cst00309) artlSeq=5036790 file=extracted/1008/5036790/2011년 제22회 감정평가사 1차시험 기출문제/1교시.hwp
%% conversion: kordoc 4.24.51 -> tools/gen1.py (auto); review=proofread-sonnet; flags=none
\begin{question}{1}
신의성실의 원칙(신의칙)에 관한 설명으로 옳지 않은 것은? (다툼이 있으면 판례에 의함)
\begin{choices}
  \item 미성년자가 법정대리인의 동의 없이 신용구매계약을 체결한 후 법정대리인의 동의 없음을 이유로 그 계약을 취소하는 것은 신의칙에 반하지 않는다.
  \item 채권자가 채권을 확보하기 위하여 제3자의 부동산을 채무자에게 명의신탁하도록 한 다음 그 부동산에 대해 강제집행을 하는 것은 신의칙에 반한다.
  \item 회사의 이사로 재직하면서 회사의 확정채무를 보증한 자는 이사직을 사임한 후에 사정변경을 이유로 그 보증계약을 해지할 수 있다.
  \item 신의성실의 원칙에 반하거나 권리남용이 되는 경우에는 당사자의 주장이 없더라도 법원은 직권으로 이를 판단할 수 있다.
  \item 징계면직처분에 불복하던 근로자가 이의 없이 퇴직금을 수령하고 다른 생업에 종사하다 징계면직일로부터 5년 후에 제기한 해고무효확인청구는 신의칙 및 실효의 원칙에 위배되어 허용될 수 없다.
\end{choices}
\end{question}
\answer{3}
